\documentclass[11pt]{article}
\usepackage[margin=1in]{geometry}
\usepackage{amsmath,amssymb,amsthm}
\usepackage{graphicx}
\usepackage[hidelinks]{hyperref}
\usepackage{microtype}

\newtheorem{theorem}{Theorem}
\newtheorem{lemma}[theorem]{Lemma}
\newcommand{\N}{\mathbb N}
\newcommand{\Z}{\mathbb Z}
\newcommand{\OB}{\mathrm{OB}}
\newcommand{\Sym}{S_\infty}

\title{Property \texorpdfstring{$(\OB)$}{(OB)} Need Not Pass to Open Subgroups}
\author{Candidate counterexample packet for arXiv:1203.6047\\Model: GPT5.6}
\date{11 August 2026}

\begin{document}
\maketitle

\noindent\textbf{Status.}
\texttt{candidate\_counterexample\_likely\_valid\_novelty\_unconfirmed}.

\section{Source problem}

Christian Rosendal asks the following in Problem 1.16 of
\emph{Global and local boundedness of Polish groups} (PDF page 14 of the
arXiv version) \cite{Rosendal2013}:

\begin{quote}
Suppose $G$ is a Polish group with property $(\OB)$ and $H\leq G$ is an
open subgroup. Does $H$ have property $(\OB)$?
\end{quote}

\begin{figure}[ht]
  \centering
  \includegraphics[width=\textwidth]{figures/open_problem_crop.png}
  \caption{The source question and its immediate context, arXiv:1203.6047,
  PDF page 14.}
\end{figure}

For a Polish group $K$, we use Rosendal's equivalent bounded-width
characterization \cite[Theorem 1.11]{Rosendal2013}:
\begin{equation}\label{eq:obcriterion}
 K\text{ has property }(\OB)
 \quad\Longleftrightarrow\quad
 \begin{minipage}{0.61\textwidth}
 for every nonempty symmetric open $V\subseteq K$, there are a finite
 $F\subseteq K$ and $k\geq1$ such that $K=(FV)^k$.
 \end{minipage}
\end{equation}
Equivalently, property $(\OB)$ fails if the group admits an unbounded
continuous length function \cite[Theorem 1.11]{Rosendal2013}.

\section{Construction and proof intuition}

Let $P=\Sym$ be the full symmetric group of $\N$, with the topology of
pointwise convergence, and let $N=\Z^{\N}$ have the product topology.  The
group $P$ acts continuously on $N$ by permuting coordinates,
\[
 (p\cdot a)_i=a_{p^{-1}(i)}.
\]
Set $G=N\rtimes P$.

The key idea is that a neighbourhood of the identity only constrains finitely
many coordinates.  A single finite permutation moves those constrained
coordinates to a disjoint block, where arbitrary values are allowed.  This
gives bounded width for the whole base group $N$.  The permutation part also
has bounded width because the pointwise stabilizer of a finite set has only
finitely many double cosets in $P$.  Together these observations verify
\eqref{eq:obcriterion} for $G$.

On the other hand, once the permutation part is required to fix coordinate
$0$, projection onto that coordinate becomes a homomorphism onto $\Z$.  The
absolute value of this homomorphism is an unbounded continuous length
function.  Thus the open coordinate stabilizer fails property $(\OB)$.

\section{Formal counterexample}

We use additive notation on $N$.  The multiplication in $G$ is
\[
 (a,p)(b,q)=(a+p\cdot b,pq).
\]

For a finite set $A\subseteq\N$, define
\[
 N_A=\{a\in N:a_i=0\text{ for all }i\in A\},\qquad
 P_A=\{p\in P:p(i)=i\text{ for all }i\in A\},
\]
and $U_A=N_A\rtimes P_A$.

\begin{lemma}\label{lem:basis}
The subgroups $U_A$, with $A\subseteq\N$ finite, form a neighbourhood basis
at the identity of $G$.
\end{lemma}

\begin{proof}
Since $P_A$ fixes every coordinate in $A$, it preserves $N_A$, so $U_A$ is a
subgroup.  A basic identity neighbourhood in the product topology constrains
only finitely many coordinates of $N$ and requires the permutation to fix
only finitely many indices.  Enlarging the union of these two finite sets to
$A$ gives $U_A$ contained in that neighbourhood.
\end{proof}

\begin{lemma}\label{lem:doublecosets}
For finite $A\subseteq\N$, the double-coset space $P_A\backslash P/P_A$ is
finite.
\end{lemma}

\begin{proof}
To $p\in P$, associate the partial injection
\[
 \theta_p:\{i\in A:p(i)\in A\}\longrightarrow A,\qquad \theta_p(i)=p(i).
\]
Left or right multiplication by an element of $P_A$ does not change
$\theta_p$.

Conversely, suppose $p,q\in P$ induce the same partial injection.  On the
indices $i\in A$ sent outside $A$, the points $p(i)$ and $q(i)$ form two
equally sized tuples of distinct points in $\N\setminus A$.  Choose
$u\in P_A$ with $u(p(i))=q(i)$ for all such $i$.  On the remaining indices
of $A$, equality of the partial injections and the fact that $u$ fixes $A$
give $up(i)=q(i)$.  Hence $(up)^{-1}q$ fixes $A$ pointwise; writing this
permutation as $v\in P_A$ gives $q=upv$.  Thus the double cosets are
classified by partial injections of the finite set $A$ into itself, of which
there are only finitely many.
\end{proof}

\begin{theorem}\label{thm:counterexample}
There is a Polish group with property $(\OB)$ having an open subgroup without
property $(\OB)$.  In particular, the answer to Rosendal's Problem 1.16 is
negative.
\end{theorem}

\begin{proof}
Both $N=\Z^{\N}$ and $P=\Sym$ are Polish, and the coordinate-permutation
action is continuous.  Hence $G=N\rtimes P$ is Polish.

We first prove that $G$ has property $(\OB)$.  Fix an identity neighbourhood
$W\subseteq G$.  By Lemma \ref{lem:basis}, choose finite $A$ with
$U_A\subseteq W$.  Choose a finite set $B\subseteq\N\setminus A$ with
$|B|=|A|$, and let $\tau\in P$ be an involution exchanging $A$ and $B$ and
fixing the remaining coordinates.

For $a\in N$, let $x\in N_A$ agree with $a$ on $\N\setminus A$ and vanish on
$A$.  Then $a-x$ is supported on $A$.  Put
$y=\tau^{-1}\cdot(a-x)=\tau\cdot(a-x)$; this is supported on $B$ and hence
belongs to $N_A$.  Directly from the multiplication law,
\begin{equation}\label{eq:basefactor}
 (a,1)=(x,1)(0,\tau)(y,1)(0,\tau).
\end{equation}

By Lemma \ref{lem:doublecosets}, choose a finite set $E\subseteq P$ of
representatives for $P_A\backslash P/P_A$.  For every $p\in P$, write
$p=uev$ with $u,v\in P_A$ and $e\in E$.  Combining this with
\eqref{eq:basefactor} shows
\[
 (a,p)\in U_A(0,\tau)U_A(0,\tau)U_A(\{0\}\times E)U_A.
\]
If
\[
 F=\{(0,1),(0,\tau)\}\cup(\{0\}\times E),
\]
then
\begin{equation}\label{eq:width}
 G=(F U_A)^4\subseteq(FW)^4.
\end{equation}

To match the exact form of \eqref{eq:obcriterion}, let $V\subseteq G$ be any
nonempty symmetric open set.  Then $V^2$ is an identity neighbourhood.  Apply
\eqref{eq:width} with $W=V^2$ to get $G=(FV^2)^4$.  After adjoining the
identity to $F$, each factor $f v_1v_2$ splits as
$(f v_1)(1v_2)$, and therefore $G=(F'V)^8$ for a finite $F'$.  Criterion
\eqref{eq:obcriterion} proves that $G$ has property $(\OB)$.

Now let $P_0=P_{\{0\}}$ and
\[
 H=N\rtimes P_0.
\]
This is an open subgroup of $G$.  Define
\[
 \varphi:H\longrightarrow\Z,\qquad \varphi(a,p)=a_0.
\]
If $p\in P_0$, then $(p\cdot b)_0=b_{p^{-1}(0)}=b_0$, so the multiplication
law gives
\[
 \varphi((a,p)(b,q))=a_0+b_0=\varphi(a,p)+\varphi(b,q).
\]
Thus $\varphi$ is a continuous surjective homomorphism.  Consequently
$\ell(a,p)=|a_0|$ is an unbounded continuous length function on $H$.
Rosendal's length-function characterization implies that $H$ does not have
property $(\OB)$.
\end{proof}

\section{Verification notes}

The proof is non-computational.  The following points were checked directly.
\begin{itemize}
\item The action $P\curvearrowright N$ is continuous in the stated Polish
topologies.
\item $U_A$ is a subgroup because $P_A$ preserves the condition of vanishing
on $A$.
\item Lemma \ref{lem:doublecosets} proves, rather than assumes, the finite
double-coset assertion.
\item Expanding the semidirect-product law in \eqref{eq:basefactor} yields
$x+\tau\cdot y=x+(a-x)=a$ and permutation part $\tau^2=1$.
\item The four blocks in \eqref{eq:width} are
$(1U_A)((0,\tau)U_A)((0,\tau)U_A)((0,e)U_A)$.
\item The map $\varphi$ is a homomorphism exactly because all permutation
parts in $H$ fix $0$.
\end{itemize}

\section{Scope and novelty status}

This construction completely answers the yes/no source problem if the proof
is accepted.  It does not claim a classification of those property-$(\OB)$
Polish groups whose open subgroups retain property $(\OB)$.

A bounded novelty search on 11 August 2026 checked the run's registry,
solution, attempt, and proof-gap indexes and searched exact source wording,
both problem labels, and combinations of ``property $(\OB)$'', ``open
subgroup'', ``counterexample'', ``wreath product'', and the displayed
semidirect product.  No later paper claiming an answer was found.  This was
not an exhaustive citation search, so novelty remains unconfirmed and should
be checked by a specialist before any public originality claim.

\section{Human-review recommendation}

High priority.  The central audit is the bounded-width factorization proving
property $(\OB)$ for $\Z^{\N}\rtimes\Sym$.  The rest of the obstruction is the
explicit quotient map onto $\Z$.

\begin{thebibliography}{9}
\bibitem{Rosendal2013}
Christian Rosendal,
\emph{Global and local boundedness of Polish groups},
Indiana Univ. Math. J. \textbf{62} (2013), no. 5,
arXiv:\href{https://arxiv.org/abs/1203.6047}{1203.6047},
\href{https://doi.org/10.1512/iumj.2013.62.5133}{doi:10.1512/iumj.2013.62.5133}.
\end{thebibliography}

\end{document}

